%%%%%%%%%%%%%%%%%%%%%%%%%%%%%%%%%%%%%%%%%%%%%%%%%%%%%%%%%%%%%%%%%%%%%%%%%%%
% Handout und Uebungsblatt zu Lehreinheit 11
% Im Skript am Kapitelende, in handouts.tex als Druckvorlage.
%%%%%%%%%%%%%%%%%%%%%%%%%%%%%%%%%%%%%%%%%%%%%%%%%%%%%%%%%%%%%%%%%%%%%%%%%%%
\begin{handout}{Wissenschaftliches Schreiben}

\textbf{Worum es heute ging.} Ein Zauber-Sprechtext ließ ein Seil
\enquote{wachsen} -- mit einem Seemann, einer \enquote{Energie des Raumes}
und dem Satz \enquote{Je mehr Sie glauben, desto länger wird es}. Beim
Übersetzen in Fachsprache blieb davon wenig übrig: geschätzte Längen,
eine offene Frage und der Hinweis, was nicht geprüft wurde. Diese
Übersetzung folgt Regeln. Wissenschaftliches Schreiben lässt sich lernen.

\textbf{Kernbegriffe}
\begin{description}
  \item[IMRaD] Introduction, Methods, Results, Discussion: Gliederung
    empirischer Berichte nach Funktion; Ergebnis und Deutung stehen getrennt.
  \item[These -- Beleg -- Schluss] Grundform des Arguments; nach Toulmin
    ergänzt um Schlussregel, Einschränkung und Ausnahme
    \parencite{toulmin1958uses}.
  \item[Paraphrase] fremder Gedanke in eigenen Worten und eigener
    Satzstruktur, \emph{mit} Beleg.
  \item[Plagiat] fremde Texte, Daten oder Ideen ohne ausreichende
    Kennzeichnung; die Absicht spielt keine Rolle.
  \item[Vierschritt] Einleitung aus Thema -- Lücke -- Frage -- Aufbau,
    nach dem Muster von \textcite{swales1990genre}.
\end{description}

\textbf{Vom Patter zur Fachsprache}

\begin{center}\small
\begin{tabular}{p{0.4\textwidth}p{0.5\textwidth}}
\toprule
\textbf{statt \ldots} & \textbf{schreiben Sie \ldots} \\
\midrule
Erzählung, Anekdote & Beobachtung mit Bedingungen \\
vage Menge (\enquote{viel}, \enquote{total}) & Wert mit Einheit und Messverfahren \\
\enquote{Studien zeigen} & die Studie, mit Beleg \\
Ursache als Gewissheit & Hypothese mit Prüfbedingung \\
\enquote{immer}, \enquote{jeder}, \enquote{sicher} & Geltungsbereich und Unsicherheit \\
Sammelbegriff (\enquote{Anker}) & Fachbegriff mit Forschungsgrundlage \\
\bottomrule
\end{tabular}
\end{center}

\textbf{Überarbeiten in zwei Runden.} Runde~1 (Inhalt): These? Belege?
Überflüssiges? Reihenfolge? Runde~2 (Sprache und Form): Füllwörter,
Nominalstil, einheitliche Begriffe, Zitierform. Erst Runde~1 abschließen.

\begin{merkbox}
Präzise schreiben heißt, bei jedem Satz angeben zu können, woher man ihn
weiß, und kenntlich zu machen, wo man es nicht weiß.
\end{merkbox}

\textbf{Zum Weiterlesen:} \textcite{apa2020manual}, Kap.~3 (Aufbau
empirischer Berichte) und Kap.~8--9 (Zitieren);
\textcite{oppenheimer2006erudite} zur Wirkung von Imponiervokabular.
\end{handout}

\begin{uebungsblatt}{Präzision und Struktur}

\aufgabe{Patter-Umbau}[in der Stunde, Kleingruppe, 10 Min.]
Den Sprechtext hören Sie gleich bei einer Vorführung; danach steht er auf
der Folie. Übersetzen Sie ihn Satz für Satz in die Sprache eines
Beobachtungsberichts. Notieren Sie zu jeder Änderung die Regel, nach der
Sie vorgegangen sind, in einem Halbsatz (etwa \enquote{Beobachtung von
Deutung trennen}).

\begin{center}
\renewcommand{\arraystretch}{2.2}
\begin{tabular}{|p{0.32\textwidth}|p{0.4\textwidth}|p{0.16\textwidth}|}
\hline
\textbf{Sprechtext (Satz)} & \textbf{Fachsprache} & \textbf{Regel} \\
\hline
 & & \\ \hline
 & & \\ \hline
 & & \\ \hline
 & & \\ \hline
\end{tabular}
\end{center}

\aufgabe{Ein schlechter Absatz, zwei Runden}[in der Stunde, Kleingruppe, 25 Min.; Abschluss als Hausaufgabe]
Der folgende Absatz ist \emph{konstruiert}, um möglichst viele Fehler zu
enthalten. Er stammt aus keiner realen Arbeit.

\begin{quote}\small
\enquote{Heutzutage weiß ja eigentlich jeder, dass Zaubertricks im
Unterricht total viel bringen. Die Schüler passen einfach viel besser auf
und merken sich alles, das sieht man ja schon, wenn man selbst mal in so
einer Stunde war. Studien haben das auch bewiesen. Manche Lehrer sind aber
dagegen, weil sie denken, es ist nur Show, was aber eigentlich nicht
stimmt, weil Staunen ja nachweislich das Gehirn aktiviert. Man sollte also
in jedem Fach zaubern. Außerdem macht es Spaß, und Spaß ist beim Lernen
das Wichtigste. Natürlich klappt nicht jeder Trick immer, aber das ist ja
egal. Ein Freund von mir ist Lehrer, und der sagt auch, dass seine Klasse
seitdem viel besser geworden ist.}
\end{quote}

\begin{enumerate}
  \item \textbf{Runde 1 -- Inhalt} (rot): Welche Thesen stecken im
    Absatz? Welche davon bräuchten einen Beleg? Was gehört nicht in einen
    wissenschaftlichen Text? Legen Sie fest, welche \emph{eine} These der
    überarbeitete Absatz vertreten soll, und ordnen Sie die Sätze neu.
  \item \textbf{Runde 2 -- Sprache und Form} (blau): Streichen Sie
    Füllwörter und Absolutheiten, vereinheitlichen Sie Begriffe,
    formulieren Sie die Unsicherheit angemessen. Setzen Sie mindestens zwei
    Belege aus der Literatur dieser Vorlesung ein.
\end{enumerate}
\schreiblinien{8}

\aufgabe{Einleitung im Vierschritt}[in der Stunde Entwurf; Hausaufgabe ca.\ 400 Wörter]
Schreiben Sie die Einleitung Ihrer Seminararbeit. Die Satzanfänge sind
Vorschläge, keine Pflicht.

\begin{center}\small
\renewcommand{\arraystretch}{1.4}
\begin{tabular}{|p{0.2\textwidth}|p{0.72\textwidth}|}
\hline
\textbf{1 Thema} & \emph{\enquote{Wenn \ldots, dann \ldots}} -- \emph{\enquote{In \ldots\ spielt \ldots\ eine zentrale Rolle, weil \ldots}} \\[8mm]
\hline
\textbf{2 Lücke} & \emph{\enquote{Untersucht wurde bisher vor allem \ldots\ (Beleg). Offen ist dagegen, ob \ldots}} \\[8mm]
\hline
\textbf{3 Frage} & \emph{\enquote{Diese Arbeit geht der Frage nach, ob/wie \ldots}} -- als echte Frage mit Fragezeichen \\[8mm]
\hline
\textbf{4 Aufbau} & \emph{\enquote{Kapitel~2 stellt \ldots\ dar. Kapitel~3 beschreibt \ldots}} \\[8mm]
\hline
\end{tabular}
\end{center}

\aufgabe{Peer-Feedback zur Einleitung}[in der Stunde, Partnerarbeit, 5 Min.]
Lesen Sie die Einleitung Ihrer Partnerin oder Ihres Partners. Beschreiben
Sie, statt zu bewerten (\enquote{Mir war nach Satz~3 nicht klar, \ldots}).

\begin{center}\small
\renewcommand{\arraystretch}{1.8}
\begin{tabular}{|p{0.5\textwidth}|c|c|c|p{0.22\textwidth}|}
\hline
\textbf{Kriterium} & \textbf{ja} & \textbf{teils} & \textbf{nein} & \textbf{Hinweis} \\
\hline
Ich weiß nach dem ersten Absatz, worum es geht. & & & & \\ \hline
Die Lücke ist benannt \emph{und} belegt. & & & & \\ \hline
Die Forschungsfrage steht als Frage da und ist beantwortbar. & & & & \\ \hline
Der Aufbau der Arbeit wird genannt. & & & & \\ \hline
Begriffe werden einheitlich verwendet. & & & & \\ \hline
\end{tabular}
\end{center}
Ein Satz, der mir hängen geblieben ist: \hrulefill\par\medskip
Eine Frage, die ich an die Arbeit habe: \hrulefill

\end{uebungsblatt}

\begin{uebungsblattb}{5}{Belegen und Nachschärfen}

\aufgabe{Plagiat, Paraphrase oder korrektes Zitat?}[Hausaufgabe]
Die Quelle ist \emph{fiktiv}: Berger, L. (2021). \emph{Staunen im Hörsaal}.
Auf Seite~42 steht:
\begin{quote}\small
\enquote{Wer staunt, bemerkt eine Lücke im eigenen Wissen. Diese Lücke
erzeugt Spannung, und Spannung sucht nach Auflösung. Lehrende können diese
Spannung nutzen, wenn sie die Auflösung nicht sofort liefern.}
\end{quote}
Ordnen Sie jede Fassung ein: \textbf{Z} korrektes Zitat, \textbf{P}
korrekte Paraphrase, \textbf{X} Plagiat, \textbf{G} Grenzfall (dann:
wovon hängt es ab?).

\begin{center}\small
\renewcommand{\arraystretch}{1.5}
\begin{tabular}{|c|p{0.7\textwidth}|p{0.1\textwidth}|}
\hline
 & \textbf{Fassung in einer Seminararbeit} & \textbf{Z/P/X/G} \\
\hline
A & \enquote{Wer staunt, bemerkt eine Lücke im eigenen Wissen} (Berger, 2021, S.~42). & \\ \hline
B & Wer staunt, bemerkt eine Lücke im eigenen Wissen. Diese Lücke erzeugt Spannung (vgl.\ Berger, 2021). & \\ \hline
C & Staunen macht eine Wissenslücke bewusst; die entstehende Spannung können Lehrende nutzen, indem sie mit der Erklärung warten (Berger, 2021, S.~42). & \\ \hline
D & Wer sich wundert, merkt eine Lücke im eigenen Wissen. Diese Lücke ruft Spannung hervor, und Spannung strebt nach Lösung. & \\ \hline
E & Dass Neugier aus einer bemerkten Informationslücke entsteht, hat Loewenstein (1994) in einer Übersichtsarbeit begründet. & \\ \hline
F & Eine Studentin übernimmt einen eigenen Absatz aus einer bereits benoteten Hausarbeit eines anderen Kurses, ohne Hinweis. & \\ \hline
G & \enquote{Das Gehirn arbeitet als Vorhersagemaschine.} -- ohne Beleg, in einer Arbeit über Predictive Processing. & \\ \hline
H & Ein Student lässt Bergers Absatz von einem KI-Werkzeug umformulieren und fügt das Ergebnis ohne Beleg ein. & \\ \hline
\end{tabular}
\end{center}

\aufgabe{Tagebuch des Staunens}[Hausaufgabe]
Suchen Sie in einem eigenen Text (Seminararbeit, Essay aus LE~1) einen
Satz, den Sie für sicher hielten und der der Frage \enquote{Woher weiß ich
das?} nicht standhält. Notieren Sie ihn, die bessere Fassung und was Sie
dabei überrascht hat.

\begin{lehrhinweis}[Erwartungshorizont]
\textbf{Aufgabe~1:} Mindestens fünf der sieben Regeln aus dem Skript
sollten auftauchen. Gute Gruppen bemerken, dass der Satz \enquote{Je mehr
Sie daran glauben \ldots} eine Immunisierung ist: Wenn das Seil nicht
wächst, hat man eben nicht genug geglaubt.

\textbf{Aufgabe~2, Runde~1:} Der Absatz enthält mindestens fünf Thesen
(Zaubern fördert Aufmerksamkeit; fördert Behalten; ist bewiesen; Staunen
aktiviert das Gehirn; Spaß ist das Wichtigste) und keinen einzigen Beleg;
die Anekdote vom Freund ist kein Beleg, der Satz über misslingende Tricks
gehört nicht zum Argument. Eine sinnvolle Entscheidung ist, nur die These
\enquote{Zaubertricks können Lernen fördern} zu vertreten und den
Gegeneinwand (\enquote{nur Show}) ernst zu nehmen. \textbf{Runde~2},
eine mögliche Endfassung:
\enquote{Ob Zaubertricks im Unterricht das Lernen fördern, ist bislang
wenig untersucht. Theoretisch ist ein Effekt plausibel: Neugier entsteht,
wenn Lernende eine Lücke in ihrem Wissen bemerken \parencite{loewenstein1994},
und in einem Experiment erinnerten Versuchspersonen die Antworten auf
Wissensfragen besser, wenn die Fragen sie zuvor neugierig gemacht hatten
\parencite{gruber2014}. Ein Zaubertrick könnte eine solche Lücke erzeugen.
Dem steht entgegen, dass eine mitreißende Darbietung die gefühlte
Lernleistung erhöhen kann, ohne dass mehr gelernt wird
\parencite{naftulin1973drfox,peerbabad2014drfox}. Ob Zaubertricks mehr
bewirken als Begeisterung, müsste daher mit Lerntests geprüft werden, nicht
mit Zufriedenheitsurteilen.}

\textbf{Aufgabe~3/4:} Häufigster Mangel ist eine unbelegte Lücke
(\enquote{Dazu gibt es noch keine Forschung}); zweithäufigster eine Frage,
die keine Frage ist (\enquote{Diese Arbeit befasst sich mit \ldots}).

\textbf{Aufgabe~5:} A~=~Z. B~=~X (wörtliche Übernahme ohne
Anführungszeichen; die Quellenangabe heilt das nicht). C~=~P. D~=~X
(verschleiernde Paraphrase ohne Beleg). E~=~kein Bezug zu Berger,
eigenständig belegt, also unproblematisch. F~=~G, in der Regel X: Ob
Selbstplagiat vorliegt, regelt die Prüfungsordnung; ohne Hinweis und
Absprache ist es unzulässig. G~=~G: In einer Arbeit über Predictive
Processing ist der Satz Kernthese einer bestimmten Theorie und braucht einen
Beleg (z.\,B.\ \textcite{clark2016}). H~=~X: Das Werkzeug ändert nichts
daran, dass Bergers Gedanke ohne Beleg übernommen wird; ob die
KI-Nutzung zusätzlich zu kennzeichnen ist, regeln Hochschule und Prüfende.
\end{lehrhinweis}
\end{uebungsblattb}
